\documentclass[10pt]{article}
\usepackage[utf8]{inputenc}
\usepackage[T1]{fontenc}
\usepackage[spanish,shorthands=off]{babel}
\usepackage{aleph-comandos}
\usepackage{moodle}
\usepackage{amssymb}
\usepackage{mathpazo}
\usepackage[margin=2cm]{geometry}

\title{Reactivos adicionales de Cálculo Diferencial}
\author{Propuesta para revisión docente}
\date{}

\begin{document}
\maketitle
Se proponen seis preguntas alternativas para MATD124. Cada pregunta tiene
cuatro opciones y una sola respuesta correcta. Se aprobaron las sustituciones
A6 por la pregunta 6, A2 por la 7 y A1 por la 9 del banco principal.
Las justificaciones para revisión
docente aparecen al final del PDF, fuera del banco exportado a Moodle.

\begin{quiz}{Calculo-Diferencial-MATD124-Adicionales-Conceptuales}

% A1. Signo de la derivada, monotonía y extremos. Dificultad estimada: básica.
\begin{multi}[single=true,shuffle=false]{calculo-MATD124-adicional-01}
Una función $f$ es continua en $[0,3]$ y diferenciable en $(0,3)$.
Se sabe que $f'(x)<0$ para $0<x<1$ y que $f'(x)>0$ para $1<x<3$.
Sin conocer una fórmula para $f$, ¿cuál de las siguientes afirmaciones
es necesariamente verdadera?
    \item[fraction=0] $f$ alcanza su máximo absoluto en $x=1$.
    \item* $f$ alcanza su mínimo absoluto en $x=1$.
    \item[fraction=0] $f(0)<f(1)$.
    \item[fraction=0] $f(3)<f(1)$.
\end{multi}

% A2. Diferencia entre monotonía estricta y no estricta. Dificultad: intermedia.
\begin{multi}[single=true,shuffle=false]{calculo-MATD124-adicional-02}
Sea $f$ diferenciable en un intervalo abierto $I$. Si $f'(x)\geq0$
para todo $x\in I$, ¿cuál de las siguientes conclusiones está garantizada?
    \item[fraction=0] Para cualesquiera $x_1,x_2\in I$ con $x_1<x_2$,
    se cumple $f(x_1)<f(x_2)$.
    \item[fraction=0] $f$ es constante en $I$.
    \item[fraction=0] $f$ no puede ser constante en ningún subintervalo de $I$.
    \item* Para cualesquiera $x_1,x_2\in I$ con $x_1<x_2$,
    se cumple $f(x_1)\leq f(x_2)$.
\end{multi}

% A3. Interpretación geométrica del valor medio. Dificultad: básica/intermedia.
\begin{multi}[single=true,shuffle=false]{calculo-MATD124-adicional-03}
Una función $f$ es continua en $[1,4]$ y diferenciable en $(1,4)$,
con $f(1)=2$ y $f(4)=8$. ¿Qué afirmación está garantizada por estas
condiciones?
    \item* Existe al menos un $c\in(1,4)$ en el que la recta tangente
    a la gráfica de $f$ es paralela a la recta que une $(1,2)$ y $(4,8)$.
    \item[fraction=0] La derivada de $f$ es igual a $2$ en todos los
    puntos de $(1,4)$.
    \item[fraction=0] $f$ es estrictamente creciente en todo $[1,4]$.
    \item[fraction=0] Existe exactamente un $c\in(1,4)$ tal que $f'(c)=2$.
\end{multi}

% A4. Papel de la diferenciabilidad en Rolle. Dificultad: intermedia.
\begin{multi}[single=true,shuffle=false]{calculo-MATD124-adicional-04}
Considere $f(x)=|x|$ en $[-1,1]$. Aunque $f(-1)=f(1)$, no existe
un $c\in(-1,1)$ tal que $f'(c)=0$. ¿Cuál de las siguientes afirmaciones
explica por qué este ejemplo no contradice el resultado que garantiza
una tangente horizontal entre extremos de igual altura?
    \item[fraction=0] La función no es continua en $x=0$.
    \item[fraction=0] El intervalo $[-1,1]$ no es cerrado.
    \item* La función no es diferenciable en $x=0$, que es un punto interior.
    \item[fraction=0] La función no tiene un mínimo en $[-1,1]$.
\end{multi}

% A5. Acotación de incrementos por la derivada. Dificultad: intermedia.
\begin{multi}[single=true,shuffle=false]{calculo-MATD124-adicional-05}
Una función $f$ es continua en $[1,3]$ y diferenciable en $(1,3)$.
Además, $f(1)=3$ y $2\leq f'(x)\leq4$ para todo $x\in(1,3)$.
¿Cuál de las siguientes afirmaciones sobre $f(3)$ está garantizada?
    \item[fraction=0] $4\leq f(3)\leq8$.
    \item* $7\leq f(3)\leq11$.
    \item[fraction=0] $5\leq f(3)\leq7$.
    \item[fraction=0] $f(3)=9$.
\end{multi}

% A6. Existencia y unicidad de un cero. Dificultad: intermedia.
\begin{multi}[single=true,shuffle=false]{calculo-MATD124-adicional-06}
Una función $f$ es continua en $[0,2]$ y diferenciable en $(0,2)$.
Se sabe que $f(0)=-1$, $f(2)=3$ y $f'(x)>0$ para todo $x\in(0,2)$.
¿Qué se puede concluir sobre las soluciones de $f(x)=0$ en $(0,2)$?
    \item[fraction=0] No existe ninguna solución.
    \item[fraction=0] Existe al menos una solución, pero las condiciones
    no permiten determinar si es única.
    \item[fraction=0] Existen al menos dos soluciones distintas.
    \item* Existe exactamente una solución.
\end{multi}

\end{quiz}

\clearpage
\section*{Clave y justificaciones para revisión docente}
\textbf{Clave:} A1: B; A2: D; A3: A; A4: C; A5: B; A6: D.

\begin{enumerate}
\item[\textbf{A1.}] \textbf{B.} El signo negativo de $f'$ implica que $f$
es estrictamente decreciente en $[0,1]$; el signo positivo implica que es
estrictamente creciente en $[1,3]$. Ambas afirmaciones se justifican
aplicando el teorema del valor medio a pares de puntos de cada intervalo.
Por tanto, $f(1)<f(x)$ para todo $x\in[0,3]$ distinto de $1$.

\item[\textbf{A2.}] \textbf{D.} Para $x_1<x_2$, el teorema del valor medio
da $f(x_2)-f(x_1)=f'(c)(x_2-x_1)\geq0$ para algún $c\in(x_1,x_2)$.
Esto garantiza que $f$ es no decreciente. No garantiza crecimiento
estricto: una función constante satisface la hipótesis. Tampoco obliga
a que $f$ sea constante, como muestra $f(x)=x$.

\item[\textbf{A3.}] \textbf{A.} La pendiente de la secante es
$(8-2)/(4-1)=2$. Por el teorema del valor medio, existe $c\in(1,4)$
con $f'(c)=2$. El resultado no exige que la derivada sea constante ni
que ese punto sea único: $f(x)=2x$ cumple las hipótesis y su derivada
vale $2$ en todos los puntos. Tampoco garantiza monotonía en todo el
intervalo. Por ejemplo, $g(x)=2x+2(x-1)(x-4)$ cumple los valores de los
extremos y $g'(x)=4x-8<0$ en $(1,2)$.

\item[\textbf{A4.}] \textbf{C.} La continuidad y la igualdad de los valores
en los extremos no bastan para aplicar el teorema de Rolle: se requiere
diferenciabilidad en todos los puntos interiores. En $0$, las derivadas
laterales son $-1$ y $1$, por lo que $f'(0)$ no existe. En los demás
puntos interiores, la derivada vale $-1$ o $1$, nunca $0$.

\item[\textbf{A5.}] \textbf{B.} Por el teorema del valor medio,
$f(3)-f(1)=f'(c)(3-1)=2f'(c)$ para algún $c\in(1,3)$. Así,
$4\leq f(3)-3\leq8$, y por tanto $7\leq f(3)\leq11$.
Las hipótesis acotan el incremento, pero no fijan un valor único.
Las funciones $f(x)=3+2(x-1)$ y $f(x)=3+4(x-1)$ alcanzan los extremos
permitidos y descartan las otras alternativas como conclusiones necesarias.

\item[\textbf{A6.}] \textbf{D.} La continuidad y el cambio de signo
garantizan al menos un cero interior. Como $f'>0$, el teorema del valor
medio implica que $f$ es estrictamente creciente y no puede tomar el
valor $0$ en dos puntos distintos. Así, el cero existe y es único.
\end{enumerate}

\section*{Sustituciones y alternativas}
\begin{itemize}
\item A1 reemplaza al reactivo original 9: evalúa extremos a
partir del signo de la derivada, sin pedir calcularla.
\item A3 puede reemplazar al original 2 o 4 si se desea reducir el peso
del cálculo de límites y evaluar el significado geométrico de la derivada.
\item A6 reemplaza al original 6: añade el análisis de unicidad
mediante monotonía a la existencia de un cero por continuidad.
\item A2 reemplaza al original 7 y evalúa la diferencia entre monotonía
estricta y no estricta. A4 y A5 siguen como alternativas para evaluar
la necesidad de las hipótesis y la estimación de cambios.
\end{itemize}
Las sustituciones aprobadas ya están incorporadas en el banco principal.
Este documento conserva las seis propuestas adicionales para consulta.

\end{document}
